%----- Acronymes -----%

\newglossaryentry{usbl}{
     type=\acronymtype,
     name={USBL},
     description={système de positionnement acoustique par ligne de base ultra courte ou \emph{Ultra Short Base Line acoustic positioning system}},
     first={USBL (système de positionnement acoustique par ligne de base ultra courte ou \emph{Ultra Short Base Line acoustic positioning system})},
     }

\newglossaryentry{slam}{
     type=\acronymtype,
     name={SLAM},
     description={Localisation et cartographie simultanées ou \emph{Simultaneous Location And Mapping}},
     first={SLAM (Localisation et cartographie simultanées ou \emph{Simultaneous Location And Mapping})},
     }
     
\newglossaryentry{ascii}{
     type=\acronymtype,
     name={ASCII},
     description={Code américain normalisé pour l'échange d'information ou \emph{American Standard Code for Information Interchange}},
     first={ASCII (Code américain normalisé pour l'échange d'information ou \emph{American Standard Code for Information Interchange})},
     }
          
     
\newglossaryentry{lbl}{
     type=\acronymtype,
     name={LBL},
     description={système de positionnement acoustique par longue ligne de base ou \emph{Long BaseLine acoustic positioning system}},
     first={LBL (système de positionnement acoustique par longue ligne de base ou \emph{Long BaseLine acoustic positioning system})},
     }
     
\newglossaryentry{auv}{
     type=\acronymtype,
     name={AUV},
     description={Véhicule sous-marin autonome ou \emph{Autonomous Underwater Vehicle}},
     first={AUV (Véhicule sous-marin autonome ou \emph{Autonomous Underwater Vehicle})},     
     }
     
\newglossaryentry{uuv}{
     type=\acronymtype,
     name={UUV},
     description={Véhicule sous-marin non-habité ou \emph{Unmanned Underwater Vehicle}},
     first={UUV (Véhicule sous-marin non-habité ou \emph{Unmanned Underwater Vehicle})},     
     }
     
\newglossaryentry{rov}{
     type=\acronymtype,
     name={ROV},
     description={Véhicule téléguidé ou \emph{Remotely Operated Vehicle}},
     first={ROV (Véhicule téléguidé ou \emph{Remotely Operated Vehicle})},     
     }
     
\newglossaryentry{hrov}{
     type=\acronymtype,
     name={HROV},
     description={Véhicule téléguidé hybride ou \emph{Hybrid Remotely Operated Vehicle}},
     first={HROV (Véhicule téléguidé hybride ou \emph{Hybrid Remotely Operated Vehicle})},     
     }
     
\newglossaryentry{gravimob}{
    type=\acronymtype,
    name={GraviMob},
    description={Système de Gravimétrie Mobile ou \emph{Mobile Gravimetry System}},
    first={GraviMob (Système de Gravimétrie Mobile ou \emph{Mobile Gravimetry System})},
    }
    
\newglossaryentry{limog}{
    type=\acronymtype,
    name={LiMo-g},
    description={Système Léger de Gravimétrie Mobile ou \emph{Light Moving Gravimetry System}},
    first={LiMo-g (Système Léger de Gravimétrie Mobile ou \emph{Light Moving Gravimetry System})},
    }
    
\newglossaryentry{ldo}{
    type=\acronymtype,
    name={LDO},
    description={Laboratoire de Domaines Océaniques},
    first={LDO (Laboratoire de Domaines Océaniques)},
    }
    
\newglossaryentry{ifremer}{
    type=\acronymtype,
    name={IFREMER},
    description={Institut Français de Recherche pour l'Exploitation de la Mer},
    first={Institut Français de Recherche pour l'Exploitation de la Mer (IFREMER)},
    }
    
\newglossaryentry{ipgp}{
    type=\acronymtype,
    name={IPGP},
    description={Institut de Physique du Globe de Paris},
    first={Institut de Physique du Globe de Paris (IPGP)},
    }

\newglossaryentry{onera}{
    type=\acronymtype,
    name={ONERA},
    description={Office National d'Études et de Recherches Aérospatiales},
    first={Office National d'Études et de Recherches Aérospatiales (ONERA)},
    }

\newglossaryentry{shom}{
    type=\acronymtype,
    name={SHOM},
    description={Service Hydrographique et Océanographique de la Marine},
    first={Service Hydrographique et Océanographique de la Marine (SHOM)},
    }    
    
\newglossaryentry{grs80}{
    type=\acronymtype,
    name={GRS80},
    description={Système de Référence Géodésique 1980 ou \emph{Geodetic Reference System 1980}},
    }
    
\newglossaryentry{svd}{
    type=\acronymtype,
    name={SVD},
    description={Décomposition en Valeurs Singulières ou \emph{Singular Value Decomposition}},
    first={décomposition en valeurs singulières ou \emph{Singular Value Decomposition} (SVD)},
    }   

\newglossaryentry{esgt}{
    type=\acronymtype,
    name={ESGT},
    description={École Supérieure des Géomètres et Topographes},
    first={École Supérieure des Géomètres et Topographes~(ESGT)},
    } 
    
\newglossaryentry{egm08}{
    type=\acronymtype,
    name={EGM2008},
    description={\emph{Earth Gravitational Model 2008}},
    first={EGM2008 (\emph{Earth Gravitational Model 2008})},
    } 

\newglossaryentry{odp}{
    type=\acronymtype,
    name={OP},
    description={Observatoire de Paris},
    first={Observatoire de Paris (OP)},
    } 

\newglossaryentry{ign}{
    type=\acronymtype,
    name={IGN},
    description={Institut National de l'information Géographique et forestière},
    first={Institut National de l'information Géographique et forestière (IGN)},
    } 

\newglossaryentry{lareg}{
    type=\acronymtype,
    name={LAREG},
    description={Laboratoire de Recherche en Géodésie},
    first={Laboratoire de Recherche en Géodésie (LAREG)},
    } 
    
\newglossaryentry{gps}{
    type=\acronymtype,
    name={GPS},
    description={Système Global de Positionnement ou \emph{Global Positioning System}},
    first={GPS (Système Global de Positionnement ou \emph{Global Positioning System})},
    } 
    
\newglossaryentry{mems}{
    type=\acronymtype,
    name={MEMS},
    description={Microsystème électromécanique ou \emph{Microelectromechanical systems}},
    first={MEMS (Microsystème électromécanique ou \emph{Microelectromechanical systems})},
    } 
%----- Définition -----%

\newglossaryentry{i-frame}{
	name={\emph{i-frame}},
	description={le \emph{i-frame} (ou \emph{inertial-frame}) désigne le repère accompagnant la Terre dans son mouvement de translation elliptique autour du Soleil},
	}    

\newglossaryentry{e-frame}{
	name={\emph{e-frame}},
	description={le \emph{e-frame} (ou \emph{earth-frame}) désigne le repère accompagnant la Terre dans son mouvement de rotation sur elle-même},
	} 

\newglossaryentry{n-frame}{
	name={\emph{n-frame}},
	description={le \emph{n-frame} (ou \emph{navigation-frame}) désigne le repère géographique local},
	first={\emph{n-frame}},
	}    
    
\newglossaryentry{b-frame}{
	name={\emph{b-frame}},
	description={le \emph{b-frame} (ou \emph{body-frame}) désigne le repère attaché au véhicule porteur},
	first={\emph{b-frame}},
	}

\newglossaryentry{alpha-frame}{
	name={$s_{\alpha}$\emph{-frame}},
	description={le $s_{\alpha}$\emph{-frame} (resp. $s_{\beta}$\emph{-frame}) (ou \emph{system-frame}) est le repère attaché à la triade accélérométrique $\alpha$ (resp. $\beta$)},
	first={$s_{\alpha}$\emph{-frame}},
	}
	
\newglossaryentry{beta-frame}{
	name={$s_{\beta}$\emph{-frame}},
	description={le $s_{\alpha}$\emph{-frame} (resp. $s_{\beta}$\emph{-frame}) (ou \emph{system-frame}) est le repère attaché à la triade accélérométrique $\alpha$ (resp. $\beta$)},
	first={$s_{\beta}$\emph{-frame}},
	}

%----- Symboles -----%     
     
\newglossaryentry{delta}{
     type=notation,
     name={\ensuremath{\delta}},
    description={Angle de cap},
    sort={delta}
     }



